% Lesson date: 2026-10-05. Material date (Europe/Moscow): 2026-10-05.
\documentclass[a4paper,12pt]{article}
\usepackage[T2A]{fontenc}
\usepackage[utf8]{inputenc}
\usepackage[russian]{babel}
\usepackage{amsmath,amssymb,helvet,icomma,geometry,microtype}
\usepackage{tikz,pgfplots,tabularx,booktabs,enumitem,parskip,fancyhdr,setspace,xcolor}
\usetikzlibrary{calc,arrows.meta,positioning,shapes.geometric}
\pgfplotsset{compat=1.18}
\usepackage[hidelinks]{hyperref}
\geometry{margin=22mm,headheight=15pt}
\renewcommand{\familydefault}{\sfdefault}
\DeclareFontFamily{T2A}{phv}{}
\DeclareFontShape{T2A}{phv}{m}{n}{<-> lass1000}{}
\DeclareFontShape{T2A}{phv}{b}{n}{<-> lasx1000}{}
\hypersetup{pageanchor=false}
\definecolor{accent}{HTML}{315F56}
\definecolor{ink}{HTML}{17211F}
\onehalfspacing
\setlength{\parindent}{0pt}
\setlength{\emergencystretch}{3em}
\setlist[enumerate]{leftmargin=*,itemsep=8pt}
\pagestyle{fancy}\fancyhf{}
\fancyhead[L]{\small Парабола и выколотые точки}
\fancyhead[R]{\small 05.10.26}
\fancyfoot[C]{\thepage}
\newcommand{\topic}[1]{\section*{\color{accent}#1}}
\newcommand{\checkitem}[1]{\(\square\)~#1\par}
\begin{document}
\begin{titlepage}
\centering
\vspace*{25mm}
{\large Математика · подготовка к ОГЭ\par}
\vspace{18mm}
{\Huge\bfseries\color{accent} Чек-лист\par}
\vspace{10mm}
{\LARGE Парабола\par и выколотые точки\par}
\vspace{12mm}
{\large От формулы к точному графику\par}
\vfill
{\large 05.10.2026\par}
\vspace{10mm}
{\large Лёвин Артём Александрович\par эксклюзивно для Ярослава Верещагина\par}
\vspace{18mm}
\begin{tikzpicture}
\draw[accent,line width=.9pt] (0,0) .. controls (3,.5) and (6,-.5) .. (10,0);
\end{tikzpicture}
\end{titlepage}

\topic{1. Короткая разминка: условие, единицы, знаки}
\checkitem{При равенстве произведения нулю рассматриваю каждый множитель.}
\[
 (5x+2)(-x-4)=0
 \quad\Longleftrightarrow\quad
 5x+2=0\ \text{или}\ -x-4=0.
\]
Корни: $-0,4$ и $-4$. Если требуется \textbf{больший корень}, ответ: $-0,4$. Перед записью ответа прочитайте последнее предложение условия.

\checkitem{Для прямого пути в прямоугольном треугольнике применяю теорему Пифагора.}
По плану катеты занимают 16 и 12 клеток; сторона клетки соответствует 2 км. Получаем 32 км и 24 км:
\[
 s=\sqrt{32^2+24^2}=\sqrt{1600}=40\ \text{км}.
\]
Для лесной дорожки выбираем её скорость: $v=15$ км/ч.
\[
 t=\frac{s}{v}=\frac{40}{15}=\frac83\ \text{ч},
 \qquad t_{\text{мин}}=60\cdot\frac83=160\ \text{мин}.
\]
\textbf{Проверка:} километры, делённые на километры в час, дают часы. Переводите результат в единицы, указанные в вопросе.

\checkitem{Различаю квадрат отрицательного числа и минус перед квадратом.}
\[
 (-2)^2=4,\qquad -2^2=-4,\qquad -(-2)^2=-4.
\]
В записи $(-2)^2$ знак входит в основание степени. В записи $-2^2$ квадрат относится только к числу 2. При подстановке отрицательного значения используйте скобки.

\clearpage
\topic{2. Опорные точки параболы}
Для $y=ax^2+bx+c$, где $a\ne0$:
\begin{itemize}[leftmargin=*,itemsep=3pt]
\item $a>0$: ветви направлены вверх; $a<0$: вниз;
\item вершина $V(x_V;y_V)$: $x_V=-\dfrac{b}{2a}$, $y_V=f(x_V)$;
\item ось симметрии: $x=x_V$;
\item пересечения с $Ox$: решите $ax^2+bx+c=0$;
\item пересечение с $Oy$: точка $(0;c)$.
\end{itemize}
По дискриминанту $D=b^2-4ac$: при $D>0$ имеются два пересечения с $Ox$, при $D=0$ одно, при $D<0$ пересечений нет.

\textbf{Пример занятия:} $y=x^2-x-6$.
\[
 x_V=\frac12=0,5,\quad y_V=(0,5)^2-0,5-6=-6,25.
\]
\[
 x^2-x-6=(x+2)(x-3).
\]
Опорные точки: $V(0,5;-6,25)$, $(-2;0)$, $(3;0)$, $(0;-6)$. Для симметричной точки $(1;-6)$ повторный расчёт не нужен.

\begin{center}
\begin{tikzpicture}
\begin{axis}[width=12cm,height=9cm,axis lines=middle,
 xmin=-3.5,xmax=4.5,ymin=-7.5,ymax=5,
 xlabel={$x$},ylabel={$y$},unit vector ratio=1 1,
 xtick={-2,0,1,3},ytick={-6,0,8},grid=major,grid style={gray!15},
 samples=250,clip=true]
\addplot[accent,thick,domain=-3.25:4.25]{x^2-x-6};
\addplot[gray,dashed] coordinates {(.5,-7.5)(.5,8)};
\addplot[only marks,mark=*,accent] coordinates {(-2,0)(3,0)(0,-6)(1,-6)(.5,-6.25)};
\node[anchor=north west,font=\small] at (axis cs:.65,-6.3) {$V$};
\end{axis}
\end{tikzpicture}
\end{center}
\checkitem{Выбираю дополнительные значения симметрично: $x_V-h$ и $x_V+h$.}
Корней может быть меньше двух. Тогда дополнительные точки особенно полезны. Подписывайте оси и сохраняйте одинаковый масштаб единицы по каждой оси.

\clearpage
\topic{3. Сокращение дроби сохраняет ограничения}
\checkitem{Сначала записываю ОДЗ по исходному знаменателю.}
\checkitem{Раскладываю трёхчлен на множители и сокращаю при выполнении ОДЗ.}
Если $x_1,x_2$ --- действительные корни, то
\[
 ax^2+bx+c=a(x-x_1)(x-x_2).
\]
Для примера занятия:
\[
 y=\frac{(x+4)(x^2+3x+2)}{x+1},\qquad x\ne-1.
\]
\[
 x^2+3x+2=(x+1)(x+2),\quad
 y=(x+4)(x+2)=x^2+6x+8,\quad x\ne-1.
\]
Вершина $(-3;-1)$; нули $-4$ и $-2$; пересечение с $Oy$: $(0;8)$. Исключённая точка находится подстановкой в \textbf{упрощённую} формулу:
\[
 (-1)^2+6(-1)+8=3.\qquad \text{Выколотая точка: }(-1;3).
\]
\begin{center}
\begin{tikzpicture}
\begin{axis}[width=12cm,height=8cm,axis lines=middle,
 xmin=-6,xmax=1,ymin=-2.3,ymax=10,
 xlabel={$x$},ylabel={$y$},unit vector ratio=1 1,
 xtick={-4,-2,-1,0},ytick={-1,0,8},grid=major,grid style={gray!15},samples=250,clip=true]
\addplot[accent,thick,domain=-6:0.3]{x^2+6*x+8};
\addplot[only marks,mark=*,accent] coordinates {(-3,-1)(-4,0)(-2,0)(0,8)};
\addplot[only marks,mark=o,mark size=3pt,very thick,accent,mark options={fill=white}] coordinates {(-1,3)};
\node[anchor=east,font=\small] at (axis cs:-1.2,4.2) {$(-1;3)$};
\end{axis}
\end{tikzpicture}
\end{center}
\textbf{Ещё один результат занятия:} $y=x^2+5x+4$, $x\ne-3$. Вершина $(-2,5;-2,25)$; нули $-4$ и $-1$; выколотая точка $(-3;-2)$.
\textbf{Самопроверка:} сокращённый множитель исчезает из формулы, но запрет на его нуль сохраняется. Выколотая точка исключена из графика.

\clearpage
\topic{Алгоритм: график после сокращения дроби}
\vspace{5mm}
\begin{center}
\begin{tikzpicture}[node distance=9mm,
 start/.style={ellipse,draw=accent,minimum width=65mm,minimum height=10mm,align=center,font=\small},
 action/.style={rectangle,rounded corners=2mm,draw=accent,text width=110mm,minimum height=12mm,align=center,font=\small},
 line/.style={-{Stealth[length=2.5mm]},thick,accent}]
\node[start] (s) {Исходная дробь};
\node[action,below=of s] (d) {Найдите нули исходного знаменателя.\\Запишите все запрещённые значения $x$.};
\node[action,below=of d] (f) {Разложите числитель и знаменатель на множители.\\Сократите общие множители на ОДЗ.};
\node[action,below=of f] (p) {Для полученного квадратного трёхчлена\\запишите $y=ax^2+bx+c$ вместе с ОДЗ.};
\node[action,below=of p] (v) {Найдите вершину и направление ветвей,\\пересечения с осями и симметричные точки.};
\node[action,below=of v] (h) {Для каждого запрещённого $x_0$ вычислите\\$y_0=ax_0^2+bx_0+c$.};
\node[action,below=of h] (g) {Постройте параболу.\\Исключите точки $(x_0;y_0)$ пустыми кружками.};
\node[action,below=of g] (c) {Проверьте подписи осей, координаты,\\знаки и сохранение исходной ОДЗ.};
\node[start,below=of c] (e) {График готов};
\foreach \a/\b in {s/d,d/f,f/p,p/v,v/h,h/g,g/c,c/e}
\draw[line] (\a)--(\b);
\end{tikzpicture}
\end{center}

\clearpage
\topic{4. Тренировка · Путь А · 10 задач}
Оформляйте вычисления и ответы. В заданиях на график подписывайте оси, вершину, пересечения и выколотые точки, если они есть.
\begin{enumerate}
\item Найдите больший корень уравнения $(4x+3)(2x-5)=0$.
\item На плане катеты прямоугольного треугольника занимают 9 и 12 клеток. Сторона клетки соответствует $0,5$ км. Велосипедист едет по прямой дорожке, совпадающей с гипотенузой, со скоростью 15 км/ч. Сколько минут займёт поездка?
\item Вычислите: а) $(-3)^2$; б) $-3^2$; в) $-(-3)^2$.
\item Разложите на множители $x^2+5x+6$ и найдите его корни.
\item Постройте график $y=x^2-4x+3$. Укажите вершину, ось симметрии, направление ветвей и пересечения с осями.
\item Постройте график $y=-x^2+2x+3$. Укажите те же характеристики.
\end{enumerate}
\vfill
\textbf{Перед следующими заданиями:} вершина проверяется подстановкой; корни --- подстановкой в уравнение. Сравнивайте значения на равном расстоянии от оси симметрии.

\clearpage
\topic{Тренировка · продолжение}
\begin{enumerate}[start=7]
\item Для $y=x^2+2x+5$ найдите вершину, ось симметрии и пересечение с $Oy$. Определите число пересечений с $Ox$. Возьмите пару симметричных точек и постройте график.
\item Запишите ОДЗ, упростите формулу и постройте график:
\[
 y=\frac{(x+5)(x^2+3x+2)}{x+1}.
\]
\item Запишите ОДЗ, упростите формулу и постройте график:
\[
 y=\frac{(x+1)(x^2+2x-3)}{x+3}.
\]
\item Запишите ОДЗ, упростите формулу и постройте график:
\[
 y=\frac{(x-2)(x^2-5x+6)}{x-3}.
\]
\end{enumerate}
\vspace{10mm}
\checkitem{Для задач 8--10 указываю вершину и пересечения с осями.}
\checkitem{Вычисляю координаты каждой выколотой точки.}
\checkitem{Сверяю ограничения с исходной дробью.}
\checkitem{Проверяю, входит ли найденная точка пересечения в ОДЗ.}

\clearpage
\topic{5. Ответы для самопроверки}
\renewcommand{\arraystretch}{1.4}
\begin{tabularx}{\linewidth}{@{}r X@{}}
\toprule
№ & Ответ \\
\midrule
1 & $2,5$. \\
2 & $30$ минут. \\
3 & а) $9$; б) $-9$; в) $-9$. \\
4 & $(x+2)(x+3)$; корни $-2$ и $-3$. \\
5 & $V(2;-1)$; ось $x=2$; ветви вверх; $Ox$: $(1;0)$, $(3;0)$; $Oy$: $(0;3)$. \\
6 & $V(1;4)$; ось $x=1$; ветви вниз; $Ox$: $(-1;0)$, $(3;0)$; $Oy$: $(0;3)$. \\
7 & $V(-1;4)$; ось $x=-1$; ветви вверх; $Ox$: пересечений нет; $Oy$: $(0;5)$. Симметричная пара: $(0;5)$ и $(-2;5)$. \\
8 & $x\ne-1$; $y=x^2+7x+10$; $V(-3,5;-2,25)$; ось $x=-3,5$; $Ox$: $(-5;0)$, $(-2;0)$; $Oy$: $(0;10)$; выколота $(-1;4)$. \\
9 & $x\ne-3$; $y=x^2-1$; $V(0;-1)$; ось $x=0$; $Ox$: $(-1;0)$, $(1;0)$; $Oy$: $(0;-1)$; выколота $(-3;8)$. \\
10 & $x\ne3$; $y=(x-2)^2$; $V(2;0)$; ось $x=2$; $Ox$: $(2;0)$; $Oy$: $(0;4)$; выколота $(3;1)$. \\
\bottomrule
\end{tabularx}\par
\vspace{6mm}
В задачах 8--10 ветви направлены вверх. Сравните свой рисунок с вычисленными координатами и отдельно проверьте пустой кружок в исключённой точке.
\end{document}
